\chapter{はじめに}\label{cha:Introduction}

ソフトウェア開発において、システム仕様書は開発の指針となる重要なドキュメントである\cite{rasinban}。
特に、顧客の曖昧な要求を具体的な要件に落とし込む過程において、システム仕様書の品質はプロジェクトの成否を左右する\cite{quality}。
しかし、システム仕様書に記述するべき内容は、組織や対象とするシステムの特性に依存する\cite{software_engineering}。
また、システム仕様書の質は、作成者の経験や知識によって左右される\cite{knowledge}。

これらの課題を解決するために、LLM(大規模言語モデル)\cite{whatllm}を用いたソフトウェア要求仕様書(本研究でのシステム仕様書に相当)の生成の自動化を目的とした先行研究がある\cite{reference}。
先行研究では、要件定義書や会議の議事録といった自然言語のドキュメントを入力として、情報を要約、抽出、整形し、あらかじめ用意した章立てに整形した情報を分類することで、ドキュメントを自動生成する手法を提案している。
先行研究では、作成するソフトウェア要求仕様書の章立てはあらかじめ人間が用意することを前提としている。

しかし、前述したように、システム仕様書に記述するべき内容は対象とするシステムの特性に依存するため、適切な章立てを用意すること自体が、高度な知識と経験を要する困難な作業である。
また、システム仕様書の骨格となる章立てに欠陥があれば、後続の文章生成がいかに精巧であっても、システム仕様書としての品質は担保されない。
したがって、システム仕様書作成の自動化においては、用意された章立てへの情報の流し込みだけでなく、システムに応じた章立てそのものを生成するアプローチが必要となる。

そこで本研究では、「章立ての作成」に着目する。
システム仕様書作成の対象となるシステムの特性に合わせた、章立ての自動生成を目的として、LLMを活用したシステム仕様書章立て生成のためのプロンプトを提案する。
システム仕様書章立て生成プロンプトは、以下の5つの構成要素で構成する。
\begin{itemize}
  \item 役割付与
  \item ベースの章立て  
  \item 開発するシステム情報
  \item システム仕様書の定義
  \item 章立て出力指示
\end{itemize}

また、システム仕様書章立て生成プロンプトには、以下の4つのプロンプトエンジニアリング\cite{prompt_engineering}の技術を適用する。
\begin{itemize}
  \item Persona Pattern
  \item Context Manager Pattern
  \item Template Pattern
  \item One-shot Prompting
\end{itemize}

本論文で作成するシステム仕様書章立て生成プロンプトを、OpenAI\cite{openai}社が提供するLLMであるGPT-5.2\cite{gpt52}を搭載したChatGPT\cite{chatgpt}に入力として与えることで、システム仕様書の章立てを生成する。


本論文の構成は以下の通りである。

第\ref{cha:Preparation}章では、本研究において必要となる前提知識を述べる。

第\ref{cha:Proposal}章では、システム仕様書章立て生成プロンプトの構成とその利用プロセスについて詳述する。

% 第\ref{cha:Application}章では、提案するプロンプトの適用例を示す。

第\ref{cha:Evaluation}章では、提案するプロンプトを用いて生成した章立ての品質を評価する実験とその結果について述べる。

第\ref{cha:Discussion}章では、実験結果の考察と今後の課題について論じる。

第\ref{cha:Conclusion}章では、本論文のまとめを述べる。

\iffalse
本テンプレートは、KatLab における修士論文執筆のための LaTeX 環境である。
このテンプレートを使用することで、Docker 環境で簡単に LaTeX 文書をコンパイルし、PDF を生成できる。

% \section{このテンプレートの使い方}

% \subsection{環境構築}
初回セットアップは以下のコマンドで行う：
\begin{verbatim}
make setup
\end{verbatim}

本テンプレートは Docker を使用しており、環境の再現性を確保している\cite{kimura-docker}。

% \subsection{PDF の生成}
以下のコマンドで paper.pdf を生成できる：
\begin{verbatim}
make
\end{verbatim}

このコマンドを実行すると、\verb|chapters/| ディレクトリ内の .tex ファイルの変更を自動で監視し、変更があれば自動的に paper.pdf を再生成する。

% \subsection{ファイル構成}
本テンプレートは以下のように構成されている：
\begin{itemize}
  \item \verb|paper.tex|: メインファイル（この内容は通常編集不要）
  \item \verb|chapters/|: 各章の .tex ファイルを配置するディレクトリ
  \item \verb|images/|: 画像ファイルを配置するディレクトリ
  \item \verb|paper.bib|: 参考文献データベース
\end{itemize}

% \section{章立ての参照}

他の章を参照する場合は、\verb|\ref{}| コマンドを使用する。
例えば、第\ref{cha:Preparation}章では基本的な \LaTeX の書き方について説明している。

% \section{本論文の構成}

以下、本論文の構成は次のとおりである。

第\ref{cha:Preparation}章では、基本的な \LaTeX の書き方について説明する。

第\ref{cha:Function}章では、図、表、数式の挿入方法について説明する。

第\ref{cha:Implementation}章では、ソースコードの挿入方法について説明する。

第\ref{cha:Indication}章では、参考文献の引用方法について説明する。

第\ref{cha:Evaluation}章では、セクションとサブセクションの使い方について説明する。

第\ref{cha:Conclusion}章では、本テンプレートのまとめを示す。
\fi
